\subsection{Evaluation Protocol}

% ===============================================================================
\subsubsection{Evaluation Metrics}

We evaluated the models using quantitative metrics for both reconstruction quality and editing performance, under a few-shot test-time setup.   

\begin{enumerate}
    \item \textbf{Fr\'echet Inception Distance (FID)}

    The Fr\'echet Inception Distance (FID) is a widely adopted metric for evaluating the quality of images generated by deep generative models, particularly in tasks such as few-shot image generation \cite{heusel2017gans, wang2023attribute}. Rather than comparing images directly at the pixel level, FID operates in a high-level feature space extracted from a pretrained Inception-v3 network \cite{heusel2017gans}. This makes it effective for capturing both the fidelity (realism) and diversity (variety) of generated samples \cite{borji2022pros}. The FID score is computed by modeling the feature distributions of real and generated images as multivariate Gaussians \cite{heusel2017gans}. The FID score is computed by modeling the feature distributions of real and generated images as multivariate Gaussians. The formula for FID is:

    \begin{equation}
        \mathrm{FID} = \|\mu_r - \mu_g\|_2^2 + \mathrm{Tr}\left( \Sigma_r + \Sigma_g - 2\left(\Sigma_r \Sigma_g\right)^{1/2} \right)
        \end{equation}
        
        Where:
        \begin{itemize}
            \item $\mu_r$, $\Sigma_r$ are the mean and covariance of the Inception features of real images,
            \item $\mu_g$, $\Sigma_g$ are the mean and covariance of the Inception features of generated images,
            \item $\|\mu_r - \mu_g\|_2^2$ represents the squared Euclidean distance between the feature means,
            \item $\mathrm{Tr}(\cdot)$ is the trace operator, which sums the diagonal elements of a matrix,
            \item $(\Sigma_r \Sigma_g)^{1/2}$ denotes the matrix square root of the product of the covariance matrices.
        \end{itemize}
        
    This formula computes the Fr\'echet (2-Wasserstein) distance between two multivariate Gaussian distributions---one fitted to real images and the other to generated images in the feature space \cite{heusel2017gans}. \textit{Lower FID is better.} Since FID measures the distance between real and generated image distributions, a lower FID score indicates that the generated images are more similar to the real images in terms of both quality and diversity. An FID of 0 implies perfect similarity. Therefore, in few-shot image generation, models that achieve lower FID scores are considered to generate more realistic and varied samples \cite{borji2022pros, wang2023attribute}.

    \item \textbf{Learned Perceptual Image Patch Similarity (LPIPS)}
    
    The Learned Perceptual Image Patch Similarity (LPIPS) is a metric used to evaluate the perceptual similarity between two images. Instead of comparing images at the pixel level, LPIPS compares deep feature activations from a pretrained neural network (such as VGG or AlexNet), which better reflect human perception~\cite{zhang2018unreasonable}. This makes LPIPS especially useful in few-shot image generation, where we want to evaluate whether generated images are visually close to target images (for fidelity), or measure how diverse the outputs are (for diversity), depending on how it is applied. The LPIPS score between two images $x$ and $x'$ is calculated as:
    \begin{multline*}
    \text{LPIPS}(x, x') = \sum_{l} \frac{1}{H_l W_l} \sum_{h=1}^{H_l} \sum_{w=1}^{W_l}\\
    \times \left\| w_l \odot \left( \hat{y}^l_{hw}(x) - \hat{y}^l_{hw}(x') \right) \right\|^2
    \end{multline*}
    where $\hat{y}^l_{hw}(x)$ is the unit-normalized feature vector of image $x$ at layer $l$ and position $(h, w)$, $w_l$ is a vector of learned weights for each channel at layer $l$, and $H_l$, $W_l$ are the height and width of the feature map at layer $l$~\cite{zhang2018unreasonable}. When interpreting LPIPS, \textit{lower values are better} in terms of similarity between generated and real images. A lower LPIPS score indicates that the generated image is more perceptually similar to the reference image. On the other hand, a higher LPIPS score may reflect less similarity (in fidelity tasks) or greater diversity (in intra-cluster evaluations of few-shot generation)~\cite{zhang2018unreasonable}.

\end{enumerate}

% ==============================================================================
\subsubsection{Baselines for Comparison}
To contextualize the performance of our proposed few-shot image generation model, we compare it against a set of strong, widely accepted baseline methods. These baselines represent different paradigms of generative modeling and are selected to provide a fair and comprehensive comparison.

\begin{itemize}    
    \item AGE \cite{AGE} -- Leverages latent space attribute editing on pretrained GANs to produce novel-class images without retraining, by disentangling class-relevant and irrelevant features.
    
    \item SAGE \cite{ding2023stableattributegroupediting} -- Builds on AGE by aggregating support images to create stable latent embeddings and inject class-consistent diversity via attribute mixing.
\end{itemize}

% =================================================================================
